\section{Discussion}
\label{sec:discussion}

The four-level validation framework provided complementary evidence that SA-generated
synthetic patients faithfully represented the real clinical population. Level~1
confirmed that marginal distributions were preserved---a necessary but insufficient
condition that MVN also satisfied by construction. Level~2 demonstrated that SA
preserved the joint cross-visit covariance structure that MVN's rank-deficient estimate
could not faithfully represent, a difference made visible in the eigenvalue spectrum
where MVN introduced spurious variance in null dimensions while SA respected the
low-rank geometry of the data. Level~3 showed that SA could amplify clinically
meaningful rare subgroups---generating 100 synthetic PCOS patients from only 3 real
ones---while preserving subgroup-specific signatures that MVN fundamentally could not
capture. Level~4 provided orthogonal mechanistic evidence: an independent ODE model,
encoding decades of coagulation biochemistry and calibrated exclusively on real
patients, confirmed that the factor-to-thrombin-generation mapping in synthetic
patients was biologically plausible, with 86--94\% cloud overlap under TF-only
conditions. Taken together, these results established that SA-generated synthetic
patients were not merely statistically similar to real patients but were consistent
with the underlying biology at multiple levels of resolution.

The ability to generate validated synthetic cohorts from very small longitudinal
datasets has immediate practical implications for maternal health research. Rare
pregnancy complications such as PCOS, preeclampsia, and gestational thrombophilia are
difficult to study because assembling large, well-characterized longitudinal cohorts
requires years of recruitment across multiple clinical sites. Our results suggested
that SA can enable hypothesis generation, power analyses, and computational modeling
studies for these understudied populations by generating synthetic cohorts that
preserve both the statistical and mechanistic properties of the real data. The
conditional generation capability demonstrated at Level~3 is particularly relevant:
rather than requiring researchers to collect additional rare patients, SA can amplify
existing small subgroups into cohorts large enough for meaningful analysis.

The key insight underlying SA's advantage over parametric alternatives is geometric.
SA generates samples on the data manifold by interpolating between stored patient
profiles via the Hopfield energy landscape, rather than fitting a parametric
distribution that must regularize or truncate high-dimensional structure. In the
$n < p$ regime that characterizes most small clinical studies, any parametric
distribution---whether Gaussian or otherwise---faces a fundamental identifiability
problem: there are more parameters to estimate than observations to estimate them
from. SA sidesteps this by operating in a PCA-reduced space where the
memory-to-dimension ratio is favorable ($K/d_{\text{PCA}} \approx 1.28$) and by
preserving the full directional structure of the data through the attention mechanism.
The eigenvalue spectrum comparison (Supplementary Fig.~S3) made this
distinction concrete: SA's spectrum faithfully reflected the rank-18 PCA structure,
while MVN's regularized spectrum introduced artificial variance in 194 additional
dimensions.

A natural concern is whether the $\beta^*$ selection procedure---originally developed
for protein sequence generation~\citep{Varner2024SA}---transfers to clinical
coagulation data. The entropy inflection method identifies the phase transition in
the Hopfield energy landscape, which is a geometric property of $K$ patterns in $d$
dimensions, not a property of what the features represent. As long as the
memory-to-dimension ratio $K/d$ is in a comparable regime, the same inflection
criterion locates the operating point that balances novelty against fidelity. For our
dataset, the empirical inflection ($\beta^* \approx 5.27$) was close to the
theoretical prediction ($\sqrt{d} \approx 4.24$), and a sensitivity analysis across
$\beta/\beta^* \in [0.1, 3.0]$ confirmed that generation quality degraded gracefully
in both directions (Supplementary Table~S1). The multiplicity weighting for
conditional generation introduced an additional consideration: achieving
$f_{\text{target}} = 0.80$ for the PCOS subgroup required $\rho \approx 26.7$,
reducing the effective pattern count to $K_{\text{eff}} \approx 4.6$
(Supplementary Table~S2). This low $K_{\text{eff}}$ explains the larger MREs observed
for PCOS features at Level~3---the energy landscape was effectively shaped by fewer
than 5 independent patterns, limiting the diversity of generated samples and pulling
means toward the population average.

The mechanistic validation at Level~4 introduced a paradigm that we believe is broadly
applicable beyond coagulation. Rather than relying solely on statistical comparisons
between real and synthetic distributions, we tested whether synthetic patients
produced biologically consistent outputs when processed by a domain-specific
computational model. This approach is available whenever a validated mechanistic model
exists for the domain of interest---pharmacokinetic models in drug development,
physiological models in critical care, or metabolic models in oncology. The
requirement is not that the mechanistic model be a perfect predictor (ours had weak
rank correlations at the individual patient level), but that synthetic patients fall
within the same model-predicted envelope as real patients. The shared systematic bias
we observed under TF+TM conditions reinforced this interpretation: both real and
synthetic patients were processed identically by the mechanistic model, producing the
same pattern of overcorrection in the protein~C pathway.

The bootstrap Mann--Whitney analysis at Level~3 revealed that 5 of 24 (21\%)
feature--condition pairs were statistically distinguishable at equal sample sizes,
warranting careful interpretation. Notably, the mean relative errors for these
features remained small (7--15\%), indicating that the central tendencies were
well-preserved. The Mann--Whitney test is sensitive to distributional shape, not
just location, and the detected differences likely reflected SA's tendency to
compress distributional tails rather than shift means. This tail compression arises
from two sources: the PCA dimensionality reduction (18 components from 216 features),
which cannot fully represent the higher-order moments of every feature, and the
direction-magnitude decomposition, which draws magnitudes from a finite empirical
distribution of $K{=}23$ norms. Features with large real-data skewness---such as
Factor~IX (coefficient of variation $\approx$~25\%) and plasminogen---were
disproportionately affected. For the PCOS subgroup, the additional pull of synthetic
means toward the population mean (due to background pattern influence in the
multiplicity-weighted energy) contributed to the plasminogen failure (MRE 14.6\%,
$p > 0.05$ in only 68\% of replicates). These results represented an honest
trade-off: SA sacrifices some tail fidelity in exchange for preserving the
manifold geometry that parametric methods cannot capture in the $n < p$ regime.

Several additional limitations warrant discussion. First, our study used a single
longitudinal dataset of $K{=}23$ patients, and further evaluation on independent
datasets and across different clinical domains is needed to establish
generalizability. Second, the PCA dimensionality reduction assumed a linear structure
in the feature space; nonlinear manifold learning methods may better capture complex
clinical relationships, though at the cost of reduced interpretability. Third, the
mechanistic validation under TF+TM conditions revealed systematic model bias,
limiting the interpretability of Level~4 results under that condition to the
observation that real and synthetic patients shared the same bias pattern. Fourth, we
did not validate synthetic patients against clinical outcomes---our validation was
restricted to statistical and mechanistic plausibility, and prospective validation
against real clinical decision-making remains an important future direction.
